\documentclass[slidestop,serif,compress]{beamer}    %slidestop=title in up-left corner (with "slidescentered" centered); compress=navigation bars as small as possible
\usepackage[T1]{fontenc}
\usepackage{lmodern}
\usepackage[ugm]{mathdesign} %use Garamond font

%gets rid of bottom navigation bars and puts slide numbers right bottom
\setbeamertemplate{footline}[frame number]{}

%alternative footer: name, institute, page number
%\setbeamertemplate{footline}
%{%
%  \leavevmode%
%  \hbox{\begin{beamercolorbox}[wd=.5\paperwidth,ht=2.5ex,dp=1.125ex,leftskip=.3cm plus1fill,rightskip=.3cm]{author in head/foot}%
%    \usebeamerfont{author in head/foot}\insertshortauthor
%  \end{beamercolorbox}%
%  \begin{beamercolorbox}[wd=.5\paperwidth,ht=2.5ex,dp=1.125ex,leftskip=.3cm,rightskip=.3cm plus1fil]{title in head/foot}%
%    \usebeamerfont{title in head/foot}\insertshortinstitute \hfill \insertframenumber \,/\,\inserttotalframenumber
%  \end{beamercolorbox}}%
%  \vskip0pt%
%}

%gets rid of navigation symbols
\setbeamertemplate{navigation symbols}{}


%%complete themes (not recommended)
%%\usetheme{Madrid}             %"Berkeley":lhs navigation bar; 1st alternative: "Madrid" without a navigation bar but with a footer including name (institution), title, date page numbering (but then \institute{Univ. ...}); 2nd alternative: defaults without all the extras; 3rd alternative "Dresden": two bottom lines)
%%"Dresden is possibly easy to adapt to UvT colors
%instead of using a complete theme inner and outer themes can be specified in isolation
%inner theme:
\useinnertheme{rounded} %actually not necessary but nice for theorem boxes etc. but enumerate listings look a bit stupid! Alternative: default = just comment it out
%outer themes:
%\useoutertheme[subsection=false,footline=authorinstitutetitlepage]{miniframes} %other options for footer see beamer user guide p. 153


%%complete color themes: (ok but inner outer styles below are better)
%%\usecolortheme{albatross}    % background is colored (in blue) in albatross; the option [overlystylish] installs a background canva which is definitely too stylish for an economics department
%%\usecolortheme{seagull}   %for greyscale (black and white)
%%default is not bad!!!

%alternative to complete color schemes one can define an inner and outer color scheme
%inner color themes
%\usecolortheme{orchid} %deep color in box(theorem etc.) background fits with outer color scheme "whale"
\usecolortheme{rose} % less aggressive slight grey shading background (in boxes/theorems etc.) fits nicely with outer color scheme \usecolortheme{seahorse}
%outer color theme
\usecolortheme{dolphin}   %alternative to seahorse: "whale" but only if inner scheme is "orchid"
%
%\usefonttheme{structurebold}       %titles, headlines, foodlines and sidebars are bold

\usepackage{amsmath, eurosym, textcomp, graphicx, amsthm, amssymb, hyperref,tikz,xcolor,sgamevar}


\setlength {\parindent}{0cm}
\newcommand {\bi}{\begin{itemize}}
\newcommand {\ei}{\end{itemize}}
\newcommand{\be}{\begin{enumerate}}
\newcommand{\ee}{\end{enumerate}}
\newcommand{\Ra}{\Rightarrow}
\newcommand{\ra}{\rightarrow}
\newcommand{\Lra}{\Leftrightarrow}
\newcommand{\del}{\partial}
\newcommand{\bea}{\begin{eqnarray*}}
\newcommand{\eea}{\end{eqnarray*}}

\renewcommand{\Re}{\mathbb{R}}

\theoremstyle{plain}
\newtheorem{prop}{Proposition}
%\newtheorem{theorem}{Theorem}
\newtheorem{assumption}{Assumption}
%\newtheorem{lemma}{Lemma}
\newtheorem{proposition}{Proposition}
\newtheorem{observation}{Observation}
%\newtheorem{corollary}{Corollary}
\theoremstyle{definition}
\newtheorem{ex}{Example}
\newtheorem{exercise}{Exercise}
%\newtheorem{definition}{Definition}

%\setlength{\parskip}{1.5ex plus 0.5ex minus 0.5ex}
\author{Christoph Schottm�ller}
\institute{University of Copenhagen}
\title{Game Theory\\Correlated equilibrium\footnote{License: CC Attribution ShareAlike 4.0}}
\date{}

\newcommand{\fr}{\frame[allowframebreaks]}
\newcommand{\ft}{\frametitle}

\newcommand{\backupbegin}{
   \newcounter{framenumberappendix}
   \setcounter{framenumberappendix}{\value{framenumber}}
}
\newcommand{\backupend}{
   \addtocounter{framenumberappendix}{-\value{framenumber}}
   \addtocounter{framenumber}{\value{framenumberappendix}}
}

\begin{document}

\begin{frame}   % Cover slide
\titlepage
\end{frame}

% \section[Outline]{}
% \frame{\ft{Outline}\tableofcontents}


\section{Correlated equilibrium}
\label{sec:corr-equil}

\fr{\ft{Correlated equilibrium}


\begin{ex}[correlated equilibrium 1]
   \begin{table}[h]
\centering
%\begin{tabular}{l|l|l}
 \begin{game}{2}{2}
       \> L \> R\\ %\hline
U   \>5,1   \> 0,0   \\
D   \>4,4   \> 1,5
\end{game}  
%\end{tabular}
\caption{correlated equilibrium}
\label{tab:correlated_eq_Aumann}
\end{table}
Find the three Nash equilibria!
\end{ex}

\framebreak

\begin{ex}[correlated equilibrium 1]
   \begin{table}[h]
\centering
%\begin{tabular}{l|l|l}
 \begin{game}{2}{2}
       \> L \> R\\ %\hline
U   \>5,1   \> 0,0   \\
D   \>4,4   \> 1,5
\end{game}  
%\end{tabular}
\caption{correlated equilibrium}
\end{table}
\end{ex}

\begin{itemize}
\item pure strategy NE give high aggregate but very unequal payoff
\item mixed strategy equilibrium gives equal but low payoff 
\end{itemize}

\framebreak

\begin{itemize}
\item Can players get equal and high payoffs?
\item flip a coin: if tails $(U,L)$, if head $(D,R)$
\item with ``unfair coins'' any payoff in the convex hull of the NE payoffs is attainable
\item can players do even better?
\item Yes:
  \begin{itemize}
  \item randomization device with three equally likely states A,B and C
  \item  P1 gets a message iff state is A
  \item P2 gets a message iff state is C
  \end{itemize}
Then the following is an equilibrium
\begin{itemize}
\item P1 plays U when he gets a message and D otherwise
\item P2 plays R when he gets a message and L otherwise
\end{itemize}
Check that no deviation is profitable!
\begin{itemize}
\item expected payoff $1/3 (5,1)+1/3(4,4)+1/3(1,5)=(3.33,3.33)$ is outside of the convex hull of the Nash payoffs
\item correlated equilibrium can lead to higher payoffs than NE
\end{itemize}
\end{itemize}


\begin{itemize}
\item interpretation correlated equilibrium:
  \begin{itemize}
\item  both players first communicate and  \emph{construct a correlation machine together}\\
each player sees output of the machine before taking action   % Then each goes to
    % a separate room and sees next the output of the correlation
    % device. Then he chooses an action. (like procurement cartel and
    % moon phase story)
  \item \emph{ impartial mediator} gives (privately!) recommendations $a_i$ to each player according to some probability distribution \\Recommendations are self-enforcing  % No player has then an incentive not to follow the recommendation, i.e. it is \emph{self enforcing}. Note that different players can have different ex post beliefs about what a third player will do.\\
  \end{itemize}
\end{itemize}

\framebreak

We now use the second interpretation:
\begin{itemize}
\item take a strategic form game $G=\langle N,(A_i),(u_i)\rangle$
\item a probability distribution $p$ over $A$ leads to the game $G^*(p)$:
  \begin{enumerate}
  \item mediator draws an action profile $a=(a_1,\dots,a_n)$ from $A$ according to probability distribution $p$
  \item mediator reveals $a_i$ to each player $i$ (but does not reveal $a_{-i}$)
  \item each player chooses an action $a_i'\in A_i$
  \item payoff for each player $i$ is $u_i(a_1',\dots,a_n')$
  \end{enumerate}
\item pure strategy for player $i$ in $G^*(p)$ is function $s_i: A_i\rightarrow A_i$ (action as function of recommendation)
\item belief of player $i$ when getting recommendation $a_i$:
  \begin{equation*}
    \label{eq:1}
    p(a_{-i}|a_i)=\frac{p(a_i,a_{-i})}{\sum_{b_{-i}\in A_{-i}}p(a_i,b_{-i})}
  \end{equation*}
  
\end{itemize}

\begin{lemma}[MSZ Thm 8.5]
  All players following the recommendation, i.e. $s_i(a_i)=a_i$ for
  all players $i$, is an equilibrium of $G^*(p)$ if and only if
  \begin{equation}\label{eq:2}
    \sum_{a_{-i}\in A_{-i}}p(a_i,a_{-i})u_i(a_i,a_{-i})\geq \sum_{a_{-i}\in A_{-i}}p(a_i,a_{-i})u_i(a_i',a_{-i})
  \end{equation}
  for all players $i$ and all actions $a_i,a_i'\in A_i$.
\end{lemma}

\framebreak

\textbf{Proof. }Expected utility of player $i$ when  $a_i$ after receiving recommendation $a_i$ is 
\begin{equation*}
  EU_i(a_i)= \sum_{a_{-i}\in A_{-i}}\frac{p(a_i,a_{-i})}{\sum_{b_{-i}\in A_{-i}}p(a_i,b_{-i})}u_i(a_i,a_{-i}).
\end{equation*}
Expected utility of playing $a_i'$ after receiving recommendation $a_i'$ is
\begin{equation*}
  EU_i(a_i')= \sum_{a_{-i}\in A_{-i}}\frac{p(a_i,a_{-i})}{\sum_{b_{-i}\in A_{-i}}p(a_i,b_{-i})}u_i(a_i',a_{-i}).
\end{equation*}
$EU_i(a_i)\geq EU_i(a_i')$ if and only if (\ref{eq:2}) holds.\qed

\vspace*{1cm}
\begin{tiny}
  For this proof, we use the convention
  $\frac{p(a_i,a_{-i})}{\sum_{b_{-i}\in A_{-i}}p(a_i,b_{-i})}=0$ if
  $p(a_i,b_{-i})=0$ for all $b_{-i}\in A_{-i}$.
\end{tiny}


\framebreak

\begin{definition}[correlated equilibrium]
  A probability distribution $p$ over $A$ is a correlated equilibrium in the strategic form game $G=\langle N,(A_i),(u_i)\rangle$ if $s_i(a_i)=a_i$ for all players $i$ is an equilibrium of $G^*(p)$. 
\end{definition}

\framebreak

\begin{prop}
  Let $\alpha^*$ be a mixed strategy equilibrium. Then the distribution $p_{\alpha^*}$ defined by
  \begin{equation*}
    p_{\alpha^*}(a_1,\dots,a_n)=\Pi_{i=1}^n\alpha_i^*(a_i)
  \end{equation*}
is a correlated equilibrium.
\end{prop}
\textbf{Proof. }

\framebreak

\begin{corollary}
 A correlated equilibrium exists in all finite games.
\end{corollary}



\framebreak


\begin{ex}[correlated equilibrium 2]
  
 \begin{table}[h]
%\centering
%\begin{tabular}{r|c|c }
%     \multicolumn{3}{c}{A}\\
 \begin{game}{2}{2}[][][A]
       \> L \> R\\ %\hline
U   \>0,1,3   \> 0,0,0   \\ %\hline
D   \>1,1,1   \> 1,0,0  
\end{game}
%\end{tabular}
\qquad
%\begin{tabular}{r|c|c }
%     \multicolumn{3}{c}{B}\\
 \begin{game}{2}{2}[][][B]
       \> L \> R\\ %\hline
U   \>2,2,2   \> 0,0,0   \\ %\hline
D   \>2,2,0   \> 2,2,2  
\end{game}
%\end{tabular}
\qquad
%\begin{tabular}{r|c|c }
%    \multicolumn{3}{c}{C} \\ 
 \begin{game}{2}{2}[][][C]
       \> L \> R\\ %\hline
U   \>0,1,0   \> 0,0,0   \\ %\hline
D   \>1,1,0   \> 1,0,3 
\end{game}
%\end{tabular}
\caption{correlated equilibrium 2}
\end{table}
\end{ex}

\begin{ex}[correlated equilibrium 2 continued]
\begin{itemize}
\item  unique NE is (D,L,A) (check!)
\item getting the $(2,2,2)$ payoff from $(U,L,B)$ or from $(D,R,B)$ would be nice for all players but P3 has an incentive to deviate (either to A or to C)
\item this example: limiting your own information can be beneficial
\item what could a mediator do to solve this problem? 
%\item mediator proposes
 %  \begin{itemize}
%   \item $(U,L,B)$ with probability 1/2
%   \item $(D,R,B)$ with probability 1/2
%   \end{itemize}
% Check that no incentives to deviate!
% Given that 3 plays B, 1 and 2 have no incentive to deviate. For 3, (U,L) and (D,R) are equally likely, hence B is a best response. Again payoffs are higher in the correlated than in the NE.
\end{itemize}
\end{ex}

\framebreak


\begin{prop}[MSZ Thm 8.9]
  Let $G=\langle N,(A_i),(u_i)\rangle$ be a strategic form game. The set of correlated equilibria of $G$ is convex. \\That is, if $p'\in \Delta A$ and $p''\in \Delta A$ are correlated equilibria, then $p=\alpha p'+(1-\alpha) p''$ is a correlated equilibrium for any $\alpha\in[0,1]$.
\end{prop}
\textbf{Proof. } 
}

\frame{\ft{Review Questions}
  \begin{itemize}
  \item What is a correlated equilibrium and how can it be interpreted?
  \item Why can correlating recommendations (sometimes) help players to achieve a higher payoff than in any Nash equilibrium?
  \item Explain why every mixed strategy Nash equilibrium can be interpreted as a correlated equilibrium.
  \end{itemize}
reading: MSZ ch. 8\\
*reading: Jann and Schottm\"uller ``Correlated equilibria in homogenous good Bertrand competition'', \emph{Journal of Mathematical Economics}, Vol. 57, March 2015, pp. 31-37
}

\fr{\ft{Exercises}
  \begin{enumerate}
  \item Determine all pure and mixed Nash equilibria of the following game. Find a correlated equilibrium in which the sum of the players' payoff is higher than in any Nash equilibrium.
   \begin{table}[h]
\centering
%\begin{tabular}{l|l|l}
 \begin{game}{2}{2}
       \> L \> R\\ %\hline
U   \>0,0   \> 6,2   \\
D   \>2,6   \> 5,5
\end{game}  
%\end{tabular}
%\caption{correlated equilibrium}
\end{table}
\item repeat the previous exercise with the following game
   \begin{table}[h]
\centering
%\begin{tabular}{l|l|l}
 \begin{game}{3}{3}
       \> L \> C \>R\\ %\hline
U   \>1,1   \> 2,4 \> 4,2   \\
M   \>4,2   \> 1,1 \> 2,4 \\
D  \> 2,4    \>4,2  \>1,1
\end{game}  
%\end{tabular}
%\caption{correlated equilibrium}
\end{table}
  \item Show that all the actions that are played with positive probability in a correlated equilibrium are rationalizable.
  \item *Think of a Bertrand game (2 firms with zero costs set price; one consumer buys from the firm with lowest price if this price is less than his valuation 1\$). Assume that prices have to be in whole cents. Show that there is no correlated equilibrium that leads to higher total payoffs than the pure strategy Nash equilibrium $(0.02,0.02)$. \\
Hint: Think of the highest recommendation given with positive probability in a correlated equilibrium.
  \end{enumerate}
}

\end{document}

















to show items one by one it is easiest to insert a "\pause" command before the \item command

hyperlinks: \hyperlink{targetname}{\beamergotobutton{Text in button}} }
    to define a target insert the option [label=targetname] after the \frame command: \frame[label=targetname]{\ft{title}...}
    alternative the command \hypertarget{targetname}{} can be used
    To go back there exists a \beamerreturnbutton{text in button}
